\subsection{Head-to-Head Comparison with Caribou (SOSP'24)}
\label{sec:evaluation:caribou}

\citet{gsteiger2024caribou} introduced Caribou, a framework for
fine-grained geospatial shifting of serverless \emph{workflows}
across geo-distributed regions. Caribou periodically (typically
hourly) re-solves an optimal deployment plan via Heuristic-Biased
Stochastic Sampling (HBSS), assigning each workflow node to a
region for the upcoming deployment window. All invocations of that
node are routed to its assigned region until the next re-deployment.
The two design points distinguishing Caribou from \GreenFaaS{} are:
\emph{(i)} decision granularity --- Caribou re-decides per-window
(hourly to daily); \GreenFaaS{} re-decides per-invocation, and
\emph{(ii)} decision scope --- Caribou's HBSS solver optimizes across
a workflow DAG; \GreenFaaS{} treats each invocation independently.

To enable a head-to-head comparison, we port Caribou to our
per-invocation simulation harness. Because the Azure 2021 trace
provides independent invocations rather than workflow DAGs, each
function is treated as a single-node ``workflow'' (a special case
explicitly supported by Caribou's DAG model). HBSS then reduces to
its degenerate form: enumerate candidate regions and pick the one
with lowest forecasted mean carbon over the upcoming window,
subject to the function's RTT budget. We give Caribou the most
favourable framing we can justify: perfect carbon forecasting over
the window (using the actual trace values rather than a Holt-Winters
estimate, eliminating one source of disadvantage); permissive
compliance constraints; zero transmission-carbon overhead (FaaS
invocations are small enough that this term is negligible); and we
test four re-deployment cadences ($1$h, $3$h, $6$h, $24$h).

\paragraph{Results on the headline topology.}
On the full 5-region topology with real Azure 2021 workload and real
EM Jan 2021 carbon, Caribou and \GreenFaaS{} produce identical
results across all four re-deployment cadences:

\begin{center}
\small
\begin{tabular}{lrr}
\toprule
Scheduler         & Carbon (g) & vs FIFO \\
\midrule
FIFO              & 231.33     & 0.00\% \\
Spatial           & 73.11      & +68.39\% \\
Caribou ($1$h)    & 73.11      & +68.39\% \\
Caribou ($3$h)    & 73.11      & +68.39\% \\
Caribou ($6$h)    & 73.11      & +68.39\% \\
Caribou ($24$h)   & 73.11      & +68.39\% \\
\GreenFaaS        & 73.22      & +68.35\% \\
\bottomrule
\end{tabular}
\end{center}

\noindent The reason for this striking convergence is that on the
5-region topology over the 14-day window, France's
nuclear-dominated grid produces the lowest carbon in \emph{every}
$1$-hour, $3$-hour, $6$-hour, and even $24$-hour bucket. Whether you
re-deploy hourly or once at the start, FR always wins, and all three
algorithms converge on the same routing. \GreenFaaS{} differs from
Spatial/Caribou by 0.04 percentage points because its
\textsc{ScoreSlots} step sometimes prefers a region with an existing
warm container over the absolutely cleanest region, trading a small
carbon-intensity penalty for cold-start avoidance; on this topology
that trade-off nets to a sub-percentage-point difference.

\paragraph{Results on a topology where granularity matters.}
The interesting comparison appears when no single region dominates
the trace. In the DE/GB/CAISO topology, GB wins approximately
$54\%$ of hours and US-CAISO $46\%$ over the 14-day trace, with the
winning region varying through the day according to solar/wind
output in California and gas-fired generation in Britain.

\begin{table}[h]
\centering\small
\caption{Caribou comparison on the time-varying topology (DE/GB/CAISO,
real Azure 2021 + real EM Jan 2021, 24h window, RTT budget 200ms).}
\label{tab:caribou_time_varying}
\begin{tabular}{lrr}
\toprule
Scheduler         & Carbon (g) & vs FIFO \\
\midrule
FIFO              & 223.90     & 0.00\% \\
Wait-Awhile       & 231.85     & $-3.55$\% \\
Lechowicz         & 235.48     & $-5.17$\% \\
Spatial           & 196.61     & +12.19\% \\
Caribou ($1$h)    & 196.64     & +12.17\% \\
Caribou ($3$h)    & 196.44     & +12.26\% \\
Caribou ($6$h)    & 198.52     & +11.34\% \\
Caribou ($24$h)   & 213.38     & +4.70\% \\
GreenFaaS-v1      & 212.33     & +5.17\% \\
\GreenFaaS        & 196.81     & +12.10\% \\
\bottomrule
\end{tabular}
\end{table}

\noindent Three observations from this table.

\textbf{(1) Caribou at hourly cadence and \GreenFaaS{} are
empirically indistinguishable} ($12.17\%$ vs $12.10\%$, gap of
$0.07$pp). When forecasting is perfect over the deployment window,
periodic re-deployment matches per-invocation routing on the
deferrable invocations that dominate the carbon footprint.
\GreenFaaS{}'s extra per-invocation re-evaluation produces no
additional carbon savings here because the lowest-carbon region does
not change \emph{within} the hourly window.

\textbf{(2) Caribou's carbon reduction degrades with coarser
cadence.} Going from $1$h to $24$h re-deployment cuts the reduction
from $12.17\%$ to $4.70\%$, a $7.5$pp loss. The $6$-hour cadence is
the regime where the cumulative within-window variation begins to
matter; the $24$-hour cadence essentially picks the daily-mean
lowest-carbon region and misses all intra-day shifts. This is a
genuine cost of Caribou's coarser-grained design point, exposed when
the optimal region varies through the day. (The $3$h cadence
marginally edges $1$h here, $+12.26\%$ vs $+12.17\%$ --- a $0.09$pp
artifact of this specific 24h window, where the $3$h boundaries
happen to align slightly better with the carbon minima than the $1$h
boundaries. The window-shift analysis below confirms the finer
cadences match or beat the coarser ones in expectation across
windows.)

\textbf{(3) Wait-Awhile and Lechowicz fail on this topology too,
losing $3$--$5\%$ to FIFO.} The failure mode of batch-oriented
carbon-aware techniques on FaaS is robust across the topologies
we test: the same conclusion from
\S\ref{sec:evaluation:lechowicz} reproduces here.

\paragraph{Window-shift stability.}
Re-running the time-varying topology over three non-overlapping 24h
windows of the Azure trace confirms the convergence is stable, not a
single-window artifact: Caribou($1$h) achieves $20.9\% \pm 18.4\%$
vs FIFO, Spatial $20.9\% \pm 19.1\%$, and \GreenFaaS{}
$21.1\% \pm 18.8\%$ --- the three track within $0.3$pp of each other
across windows. Caribou($24$h) trails at $15.8\% \pm 23.5\%$, a
$5$pp gap that persists across windows. The large absolute standard
deviations reflect the day-to-day variation in workload intensity
(the same effect characterized in \S\ref{sec:evaluation:realboth});
the relative ordering between schedulers is what remains stable.

\paragraph{Summary of the head-to-head.}
The Caribou comparison supports a precise, honest statement of
\GreenFaaS{}'s contribution: \emph{per-invocation routing does not
improve over hourly periodic re-deployment on the workloads and
carbon traces we tested, when forecasting is perfect.} The
contribution of \GreenFaaS{} over Caribou is not in carbon savings
on the headline topology, but in two other dimensions:
\emph{(i)} no requirement for workflow-DAG specification (Caribou
requires the developer to annotate workflow structure;
\GreenFaaS{} operates on independent invocations); and
\emph{(ii)} no requirement for periodic re-deployment infrastructure
(Caribou requires a per-region Knative-like control plane to enact
re-deployments; \GreenFaaS{} requires only a per-invocation
proxy). When the developer cannot or does not want to annotate
DAGs, or when the deployment infrastructure does not support
controlled re-deployment cycles, \GreenFaaS{} provides the same
carbon-reduction benefit through a simpler operational surface.

For topologies with a persistently dominant low-carbon region (such
as the France-included 5-region topology, where FR wins every hour of
the 14-day trace), Caribou and \GreenFaaS{} converge to the same
routing on the configurations we tested. For topologies where the
optimal region varies within the day (such as DE/GB/CAISO), hourly
Caribou and per-invocation \GreenFaaS{} are both effective, while
daily Caribou leaves substantial carbon on the table. We caution that
these convergence findings rest on the specific topologies and
windows tested and on perfect within-window forecasting; we treat
them as a clarification rather than a refutation of Caribou's
contribution: their work targeted workflow-level optimization on AWS,
ours targets per-invocation on multi-region serverless platforms; the
empirical results show the two approaches converge on the carbon
dimension for our workload class.
